% !PRINTED 654
\weberpage{654}

\begin{center}
\S.\,181.\\
\gesperrt{Preparation for the Proof.}
\end{center}

After the expositions of the preceding section, all that is still required for the proof of the great Theorem I is the demonstration that the number contained in $\Omega_m$
\[
\omega = (r, x_0)^m
\]
is the $m$th power of a cyclotomic number.

1. \gesperrt{The essential property of the numbers $\omega_n$ conjugate to $\omega$, on which the proof rests, is that}
\weq{1}{\omega_n^{-1} \omega_1^n = \alpha^m}
\gesperrt{is the $m$th power of a number in $\Omega_m$.}

Here $n$ may be any number not divisible by $q$. If we form the norms of the right and left sides of (1) in the field $\Omega_m$, then, denoting the norm of the number $\alpha$, which is a rational number, by $a$, we obtain, since every $\omega_n$ has the same norm as $\omega$,
\[
N(\omega)^{n-1} = a^m.
\]
Now if $m$ is odd, we may here take $n = 2$, and obtain as a \gesperrt{conclusion from 1.}:

2. \gesperrt{The norm of $\omega$ is the $m$th power of a rational number}
\weq{2}{N(\omega) = a^m.}

% [sic] The German source prints the anomalous punctuation ``Potenz von 2.,''; the translation normalizes it without concealing the source reading.
If $m$ is a power of 2, this conclusion is no longer admissible. It then follows from 1. only that $N(\omega)^2$ is the $m$th power of a rational number. Nevertheless the numbers $\omega$, as is to be seen from (10) of the preceding section, must in this case too still satisfy condition 2., which then, however, is no longer of itself already contained in 1., but is added as a new property.

\gesperrt{What we have to prove is that from 1. and 2. it follows that $\omega$ itself is also the $m$th power of a cyclotomic number, though it may be in a higher field.}

For this, however, it is necessary to investigate the decomposition of the numbers $\omega$ into their prime factors in the field $\Omega_m$.

The prime number $q$ is, as demonstrated in \S.\,169, with the $\varphi(m)$th power of a prime number $\sigma = 1 - r$ existing in $\Omega_m$
%CONT>>
